\section{Preliminaries and Problem Statement}
\label{sec:problem}

\subsection{Evidence graph}
An \emph{evidence graph} is a directed, typed multigraph $G=(V,E,\tau_V,\tau_E,\phi)$. Each vertex $v\in V$
is an entity (a process, file, registry key, network connection, user or host) with type
$\tau_V(v)$ and properties $\phi(v)$, a map from field names to values (for a process: image path, command line,
user, start time, hashes; for a file: path, signature status). Each edge $e=(u,w)\in E$ is a relation with type
$\tau_E(e)$ (process creation, file write, file execution, registry write, network connect, process access,
\dots) and properties $\phi(e)$ (timestamp, access mask, value written). Entity and relation identifiers are
opaque strings. Windows and Linux use distinct type names where semantics differ (for example
\tool{win\_process\_create} and \tool{linux\_process\_create}).

\subsection{Alert and tools}
An \emph{alert} $a=(t, s, T_V, T_E)$ has a title $t$, a summary $s$, trigger entities $T_V\subseteq V$ and
optionally trigger relations $T_E\subseteq E$. The agent never sees $G$ directly. It interacts through a
tool set $\Omega$ advertised by the runtime with each task:
\begin{itemize}
\item $\tool{get\_entity}(v)$ returns $\phi(v)$ and the identifiers of the relations incident to $v$;
\item $\tool{get\_relation}(e)$ returns $\phi(e)$ and its endpoints;
\item $\tool{search}(q,\text{scope},\text{window})$ returns a page of entities and relations whose
identifier or properties contain the substring $q$, optionally restricted to a type and a time window;
\item $\tool{finish\_investigation}(\sigma)$ submits a final answer $\sigma$ and ends the task.
\end{itemize}
Every call except a submission may carry a short free-text field describing what it checks. A
\emph{submission} is either malicious, $\sigma=(\textsf{mal}, A, r)$ with response artifacts
$A=\{(v_i,k_i)\}$ pairing an entity with a response kind $k_i$ (for example host to isolate, identity to
rotate, persistence to remove, file to delete, network block, or an \emph{observed} kind for entities that
need no action) and a summary $r$; or benign, $\sigma=(\textsf{ben}, L, r)$ with legitimacy evidence
$L=\{(x_j,F_j)\}$ citing an entity or relation $x_j$ and the property fields $F_j$ that prove it legitimate.

\subsection{Scoring}
Let $y\in\{\textsf{mal},\textsf{ben}\}$ be the true verdict. For malicious tasks the ground truth contains a
\emph{core} set $C\subseteq V$ that a complete response must name and an \emph{acceptable} set
$C\subseteq A^\star\subseteq V\cup E$ of entities involved in the attack, which may be named without
penalty. For benign tasks it contains anchors $\mathcal{L}$, each with the set $D(x)$ of its decisive
property fields, and a \emph{neutral} set $N$ (the trigger's process lineage up to the highest anchor, its
host and accounts) whose citation neither earns nor loses credit. Quality $Q\in[0,1]$ is
\begin{equation}
Q(\sigma)=\begin{cases}0 & \text{if the verdict is wrong,}\\ \tfrac12+\tfrac12\,S(\sigma) & \text{otherwise,}\end{cases}
\end{equation}
where for malicious submissions $S$ is the F1 score of
\begin{equation}
\text{rec}=\frac{1}{|C|}\sum_{v\in C}\kappa(v),\qquad
\text{prec}=\frac{|\{v: (v,\cdot)\in A\}\cap A^\star|}{|\{v: (v,\cdot)\in A\}|},
\end{equation}
with $\kappa(v)=1$ if $v$ is named with a response kind that fits its type, $0.5$ if named with another
kind, and $0$ if not named. For benign submissions,
$S=\max\bigl(0,\ \frac{|\{x\in\mathcal L:\ \exists (x,F)\in L,\ F\cap D(x)\neq\emptyset\}|}{|\mathcal L|}-\tfrac12\pi\bigr)$,
where $\pi=o/(h+o)$, with $h$ the number of anchors credited and $o$ the number of cited items outside
$\mathcal L\cup N$. With $n$ tool calls and a task-specific
optimum $n^\star$, efficiency and reward are
\begin{equation}
E=\min\!\left(1,\frac{n^\star+1}{n}\right),\qquad R=Q\,(0.8+0.2\,E).
\end{equation}
The scorer is ours; it follows the structure of the competition's published scoring (verdict and
artifacts, penalised by wasted calls), but its exact weights are not the competition's.

\subsection{Problem statement}
Given an alert $a$ and the tools $\Omega$, an investigation agent is a policy that issues a sequence of tool
calls and one submission. The objective is to maximise expected reward $\mathbb{E}[R]$ over a distribution of
alerts, that is, to reach the correct verdict, name exactly the entities that need a response (or prove
legitimacy with decisive evidence), and do so with few calls. Calls refused by the runtime, malformed calls
and calls past the budget all count against efficiency or quality, which motivates resolving avoidable mistakes
\emph{before} they reach the runtime.

\paragraph{Assumptions} Telemetry is trustworthy (we do not model an adversary tampering with sensors), the
agent sees only tool results, and ground truth is hidden from it. Telemetry fields that contain
attacker-controlled text, such as command lines, are a potential prompt-injection channel; we discuss this in
\S\ref{sec:limitations}.

\subsection{Research questions}
\begin{itemize}
\item[\textbf{RQ1}] \emph{Effectiveness.} How accurate are the agent's verdicts and responses, and how do
false-positive and false-negative rates compare?
\item[\textbf{RQ2}] \emph{Efficiency and cost.} How many tool calls, how much time and how many tokens does an
investigation take, and where are calls wasted?
\item[\textbf{RQ3}] \emph{Ablation.} How much does each local safeguard and the investigation method
contribute to quality, efficiency and reward?
\item[\textbf{RQ4}] \emph{Robustness.} Does performance hold across operating systems and on benign
look-alikes that share an alert with an attack?
\item[\textbf{RQ5}] \emph{Shortcuts and evasion.} Does the agent rely on label cues in the recordings or on
attacker-chosen names, and how does it fare when they are removed or renamed without changing behaviour?
\end{itemize}
